\begin{abstract}
On November 20, 2026, Levent Alp\"oge announced a proof of the Riemann hypothesis that he and Claude had found together. He released the proof in Lean. Shortly thereafter, the Lean and Mathlib communities certified the formal statement and proof.

Alp\"oge's argument is as follows. First, a quantitative form of Landau's theorem on Mellin transforms reduces the zero-free half-plane $\operatorname{Re}s>1/2$ for Dirichlet $L$-functions, including the Riemann zeta function, to a one-sided square-root bound for smoothed prime sums on suitable intervals. Alp\"oge begins with this reduction and then adapts Stepanov's point-counting method over finite fields. In this arithmetic version of Stepanov's method, the increasing divisibility of successive theta lifts at inert primes replaces high-order vanishing at rational points. Each unit of $p$-adic divisibility contributes $\log p$ to the gain term in the slope inequality. For a sum of length $X$, a chain of about $\sqrt X$ theta lifts balances the errors in that inequality and gives the one-sided square-root bound that Landau's criterion requires.

This manuscript presents a human-readable reconstruction of Alp\"oge's argument. It gives the argument in conventional mathematical language and identifies the mathematical inputs and manuscript-to-Lean correspondences that still require verification.
\end{abstract}
